\documentclass[10pt,a4paper]{amsart}
\usepackage{amsmath,amssymb,amsthm,graphicx}
% Retained preamble for the Stacks Project chapter snapshot in this directory.
% The chapter's own preamble.tex is a copy of the shared Stacks Project preamble; the block
% below is sliced MECHANICALLY, line for line, from it. Nothing was added, removed or
% redefined; the wrapper class, the xr-hyper/hyperref wiring and the externaldocument list
% are omitted because the assembler supplies the document class.
%
% lines taken from preamble.tex: 3, 5, 9, 11, 12, 156, 157, 158, ... (51 lines total)
%
% package coverage restored from the chapter's own preamble (2026-10-13): xy, xy-2cell

\usepackage[all]{xy}
\xyoption{2cell}
\UseAllTwocells

\newenvironment{reference}{\comment}{\endcomment}
\newenvironment{slogan}{\comment}{\endcomment}
\newenvironment{history}{\comment}{\endcomment}
\theoremstyle{plain}
\newtheorem{theorem}[subsection]{Theorem}
\newtheorem{proposition}[subsection]{Proposition}
\newtheorem{lemma}[subsection]{Lemma}
\theoremstyle{definition}
\newtheorem{definition}[subsection]{Definition}
\newtheorem{example}[subsection]{Example}
\newtheorem{exercise}[subsection]{Exercise}
\newtheorem{situation}[subsection]{Situation}
\theoremstyle{remark}
\newtheorem{remark}[subsection]{Remark}
\newtheorem{remarks}[subsection]{Remarks}
\numberwithin{equation}{subsection}
\def\lim{\mathop{\mathrm{lim}}\nolimits}
\def\colim{\mathop{\mathrm{colim}}\nolimits}
\def\Spec{\mathop{\mathrm{Spec}}}
\def\Hom{\mathop{\mathrm{Hom}}\nolimits}
\def\Ext{\mathop{\mathrm{Ext}}\nolimits}
\def\SheafHom{\mathop{\mathcal{H}\!\mathit{om}}\nolimits}
\def\SheafExt{\mathop{\mathcal{E}\!\mathit{xt}}\nolimits}
\def\Sch{\mathit{Sch}}
\def\Mor{\mathop{\mathrm{Mor}}\nolimits}
\def\Ob{\mathop{\mathrm{Ob}}\nolimits}
\def\Sh{\mathop{\mathit{Sh}}\nolimits}
\def\NL{\mathop{N\!L}\nolimits}
\def\CH{\mathop{\mathrm{CH}}\nolimits}
\def\proetale{{pro\text{-}\acute{e}tale}}
\def\etale{{\acute{e}tale}}
\def\QCoh{\mathit{QCoh}}
\def\Ker{\mathop{\mathrm{Ker}}}
\def\Im{\mathop{\mathrm{Im}}}
\def\Coker{\mathop{\mathrm{Coker}}}
\def\Coim{\mathop{\mathrm{Coim}}}
\DeclareMathSymbol{\boxtimes}{\mathbin}{AMSa}{"02}
\def\QCohstack{\mathcal{QC}\!\mathit{oh}}
\def\Cohstack{\mathcal{C}\!\mathit{oh}}
\def\Spacesstack{\mathcal{S}\!\mathit{paces}}
\def\Quotfunctor{\mathrm{Quot}}
\def\Hilbfunctor{\mathrm{Hilb}}
\def\Curvesstack{\mathcal{C}\!\mathit{urves}}
\def\Polarizedstack{\mathcal{P}\!\mathit{olarized}}
\def\Complexesstack{\mathcal{C}\!\mathit{omplexes}}
\def\Pic{\mathop{\mathrm{Pic}}\nolimits}
\def\Picardstack{\mathcal{P}\!\mathit{ic}}
\def\Picardfunctor{\mathrm{Pic}}
\def\Deformationcategory{\mathcal{D}\!\mathit{ef}}
\usepackage{cleveref}
\crefname{originallemma}{original lemma}{original lemmas}
\Crefname{originallemma}{Original lemma}{Original lemmas}
\makeatletter
\newlabel{lemma-add-simplices}{{01GN}{source}{}{sourcepointer-1}{}}
\newlabel{lemma-add-simplices@cref}{{[originallemma][0][]01GN}{[1][0][]source}{}{}{}}
\makeatother
\crefname{originallemma}{original lemma}{original lemmas}
\Crefname{originallemma}{Original lemma}{Original lemmas}
\makeatletter
\newlabel{lemma-funny-gamma}{{01GH}{source}{}{sourcepointer-2}{}}
\newlabel{lemma-funny-gamma@cref}{{[originallemma][0][]01GH}{[1][0][]source}{}{}{}}
\makeatother
\crefname{originallemma}{original lemma}{original lemmas}
\Crefname{originallemma}{Original lemma}{Original lemmas}
\makeatletter
\newlabel{lemma-hom-hypercovering}{{01GP}{source}{}{sourcepointer-3}{}}
\newlabel{lemma-hom-hypercovering@cref}{{[originallemma][0][]01GP}{[1][0][]source}{}{}{}}
\makeatother
\title{Making two maps of hypercoverings homotopic after refinement}
\author{The Stacks Project Authors}
\date{stacks-project:hypercovering.tex:local-snapshot-20261012; GNU FDL 1.2 or later}
\begin{document}
\maketitle
\noindent\textit{Editorial note.} Statement span is lines 2074-2083 of sources/stacks-hypercovering/pinned-original-source.tex; proof span is lines 2085-2122, the complete author proof (36 body lines).\par\medskip
\noindent\textit{Editorial note.} Navigation stubs: 3 cited-but-not-included labels are supplied as non-load-bearing pointers with their pinned tags where the snapshot tag index has them.\par\medskip
\noindent\textit{Editorial note.} extra\_packages is empty on purpose: the retained preamble already loads xy with options and a second declaration clashes in the assembler wrapper.\par\medskip
\noindent\textit{Editorial note.} Text is retained verbatim, including the author's own typos and placeholders.\par\medskip
\noindent\textit{Editorial note.} Queue index 14 of queue\_hypercovering.json (pb=1390, stmt 2074-2083, proof 2085-2122).\par\medskip
\noindent\textit{Editorial note.} stacks-hypercovering wave-1 (hc01): fresh adjudication accept (2026-10-14); pool per standing fallback order (hypercovering -> pic -> spaces-more-cohomology) after md01 closed the moduli pool; hypercovering.tex 53 labels/34 with proof judged 26 accept / 8 reject / 14 pb=0 env-dup rows excluded; spans cite-free and reference/slogan/history-free so no patchenv needed.\par\medskip
\noindent\textit{Allowed standard prerequisites (editorial).}\begin{itemize}
\item Explicit prerequisite boundary. The cited inputs are prerequisites; the argument in the retained span is not replaced by any cited result.
\end{itemize}\medskip
\ifdefined\hypertarget\hypertarget{sourcepointer-1}{}\fi
\noindent\textit{Source reference.} Original lemma 01GN: sources/stacks-hypercovering/pinned-original-source.tex, source lines 1890--1904. Cited by the retained statement or proof but not part of the collected spans; supplied only as a navigation pointer. The argument does not depend on its content.\par\medskip
\ifdefined\hypertarget\hypertarget{sourcepointer-2}{}\fi
\noindent\textit{Source reference.} Original lemma 01GH: sources/stacks-hypercovering/pinned-original-source.tex, source lines 1511--1536. Cited by the retained statement or proof but not part of the collected spans; supplied only as a navigation pointer. The argument does not depend on its content.\par\medskip
\ifdefined\hypertarget\hypertarget{sourcepointer-3}{}\fi
\noindent\textit{Source reference.} Original lemma 01GP: sources/stacks-hypercovering/pinned-original-source.tex, source lines 1996--2009. Cited by the retained statement or proof but not part of the collected spans; supplied only as a navigation pointer. The argument does not depend on its content.\par\medskip
\begin{lemma}
\label{lemma-homotopy}
Let $\mathcal{C}$ be a site with fibre products.
Let $X$ be an object of $\mathcal{C}$.
Let $K, L$ be hypercoverings of $X$.
Let $a, b : K \to L$ be morphisms of hypercoverings.
There exists a morphism of hypercoverings
$c : K' \to K$ such that $a \circ c$ is homotopic
to $b \circ c$.
\end{lemma}

\begin{proof}
Consider the following commutative diagram
$$
\xymatrix{
K' \ar@{=}[r]^-{def} \ar[rd]_c &
K \times_{(L \times L)} \Hom(\Delta[1], L)
\ar[r] \ar[d] & \Hom(\Delta[1], L) \ar[d] \\
& K \ar[r]^{(a, b)} & L \times L
}
$$
By the functorial property of $\Hom(\Delta[1], L)$
the composition of the horizontal morphisms
corresponds to a morphism $K' \times \Delta[1] \to L$ which
defines a homotopy between $c \circ a$ and $c \circ b$.
Thus if we can show that $K'$ is a
hypercovering of $X$, then we obtain the lemma.
To see this we will apply Lemma \ref{lemma-funny-gamma}
to the pair of morphisms $K \to L \times L$
and $\Hom(\Delta[1], L) \to L \times L$.
Condition (1) of Lemma \ref{lemma-funny-gamma} is satisfied.
Condition (2) of Lemma \ref{lemma-funny-gamma} is true because
$\Hom(\Delta[1], L)_0 = L_1$, and the morphism
$(d^1_0, d^1_1) : L_1 \to L_0 \times L_0$ is a
covering of $\text{SR}(\mathcal{C}, X)$ by our
assumption that $L$ is a hypercovering.
To prove condition (3) of Lemma \ref{lemma-funny-gamma}
we use Lemma \ref{lemma-hom-hypercovering} above. According
to this lemma the morphism $\gamma$ of condition (3) of Lemma
\ref{lemma-funny-gamma} is the morphism
$$
\Hom(\Delta[1] \times \Delta[n + 1], L)_0
\longrightarrow
\Hom(U, L)_0
$$
where $U \subset \Delta[1] \times \Delta[n + 1]$.
According to Lemma \ref{lemma-add-simplices}
this is a covering and hence the claim has been proven.
\end{proof}

\end{document}
